\chapter{Présentation de l'entreprise}

\section{VideoStitch}
\subsection{Présentation générale}
\begin{wrapfigure}{r}{35mm}
    \includegraphics{images/videostitch.png}
    \caption{Logo de VideoStitch}
\end{wrapfigure}
VideoStitch est une entreprise française d'édition de logiciels, qui conçoit et 
vend des logiciels de capture, de montage, de montage, de diffusion et de lecture 
de vidéos panoramiques à 360\degree. Elle s'intéresse par extension au marché de 
la réalité virtuelle.
Son site internet est à l'adresse suivante : \url{http://www.video-stitch.com/}.\\
C'est une Société à Actions Simplifiées (SAS) basée sur Paris, qui emploie moins
d'une dizaine de personnes actuellement. La plupart ayant été embauchés il y a 
quelques mois, c'est une société en forte croissance salariale. 
Nicolas Burtey en est le président.

\subsection{Historique}
VideoStitch a été fondée par ce même Nicolas Burtey en janvier 2012.\\
Photographe diplômé de l'Ecole Lumière, s'étant fait connaître par la photo panoramique 
et la photo fish-eye, M. Burtey évolue par la suite dans le domaine de la vidéo panoramique.\\
\begin{wrapfigure}{l}{50mm}
    \includegraphics[width=50mm]{images/loopin.png}
    \caption{Logo de Loop'in}
\end{wrapfigure}
Il crée sa société de production vidéo, Loop'In, avec laquelle il réalise une 
certain nombre de vidéos panoramiques à 360\degree a but marketing pour le compte de clients. 
Des exemples de ces vidéos sont disponibles à cette adresse : \url{http://www.loop-in.com/galerie/}.
La vidéo panoramique, étant alors un procédé très jeune encore inconnu du grand public, est 
lente et difficile à réaliser: beaucoup de tâches doivent être faites de manière 
manuelle et répétitive, avec des outils inadaptés car non prévus à cette utilisation.
Si la création de photos panoramiques et à 360\degree est aisée, il manque toutefois 
un logiciel complet et mature qui permetrait l'édition et le montage de vidéos paramiques.\\
\newline
M. Burtey décide alors de créer VideoStitch pour concevoir cet outil, un logiciel 
d'édition et de montage de vidéos panoramiques : VideoStitch Studio.\\
Quand ce logiciel atteint une version stable, en 2014, VideoStitch oriente alors 
ses efforts vers un nouveau logiciel, Vahana VR, permettant cette fois-ci une réalisation 
et une diffusion en direct de vidéos panoramiques à 360\degree.\\
\newline
Le but de la société est de permettre une réalisation aisée et de qualité de vidéos 
panoramiques à 360\degree, de la capture des images à la diffusion de la vidéo finale, 
en passant par sa réalisation complète.


\section{Les marchés et les produits}
fake ref \cite{what-is-equirectangular}
% La vidéo 360 panoramique et la VR, les principaux acteurs, bref historique et portrait actuel

\section{Produits et stratégies de l'entreprise}
% Studio, Vahana VR
% Player, réseaux sociaux, salons, rencontres entreprises

\section{Organisation de l'entreprise}
% Organisation de l'entreprise : les différentes personnes et articulation des rôles et des tâches.
% Organisation du travail : Agile/Scrum, JIRA, Confluence, Github et autres outils.

\section{Stratégies pour l'avenir et ambitions}
% Analyse de l'avenir des marchés, stratégies de positionnement et d'investissement, ambitions
